%%% DON'T MODIFY THE PACKAGE IMPORTS & CUSTOMIZATIONS UNLESS YOU REALLY NEED TO ADD ANY PACKAGE %%%
%--------------------------------------------------
% PACKAGES %
%--------------------------------------------------
\documentclass{article}
\usepackage[T1]{fontenc}
\usepackage[utf8]{inputenc}
\usepackage{lmodern}
\usepackage{graphicx}
\usepackage{subcaption}
\usepackage{amsmath}
\usepackage{fontawesome5}
\usepackage{csquotes}
\usepackage{comment}
\usepackage{enumitem}
\usepackage[table,xcdraw]{xcolor}
\usepackage{listings}
\usepackage{authblk}
\usepackage{lipsum}
\usepackage{minitoc}
\usepackage{booktabs}
\usepackage{longtable}
\usepackage{array}
\usepackage[justification=justified]{caption}
% \captionsetup{width=\textwidth}
\usepackage{orcidlink}
\usepackage{placeins}
\usepackage{bibunits}
\usepackage{hyperref}
\usepackage{siunitx}
\usepackage{float}
\usepackage{tocloft}
\usepackage{wrapfig}

\usepackage{amsmath,amssymb,amsthm}
\usepackage{bm}
\usepackage{algorithm}
\usepackage{algpseudocode}
\usepackage{booktabs}
\usepackage{multirow}
\usepackage{tgadventor}

\usepackage[margin=1in]{geometry}
\usepackage[most]{tcolorbox}

\usepackage[version=3]{mhchem}

%--------------------------------------------------
% COLORS & CODE STYLE
%--------------------------------------------------
\definecolor{myblue}{rgb}{0.1, 0.2, 0.6}
\definecolor{mytheme}{RGB}{242, 246, 255}
\definecolor{mygray}{gray}{0.35}
\definecolor{mygreen}{RGB}{0,120,0}
\definecolor{codebg}{RGB}{248,248,248}
\definecolor{keywordcolor}{RGB}{0,0,180}
\definecolor{stringcolor}{RGB}{153,0,0}
\definecolor{commentcolor}{RGB}{0,128,0}

\lstdefinestyle{bandgap}{
  backgroundcolor=\color{codebg},
  basicstyle=\ttfamily\footnotesize,
  keywordstyle=\color{keywordcolor}\bfseries,
  stringstyle=\color{stringcolor},
  commentstyle=\color{commentcolor}\itshape,
  showstringspaces=false,
  columns=fullflexible,
  keepspaces=true,
  frame=single,
  rulecolor=\color{mygray},
  breaklines=true,
  postbreak=\mbox{\textcolor{mygray}{$\hookrightarrow$}\space},
  language=Python
}

%--------------------------------------------------
% HYPERREF
%--------------------------------------------------
\hypersetup{
    colorlinks=true,
    linkcolor=[RGB]{0,51,153},
    urlcolor=[RGB]{0,51,153},
    citecolor=[RGB]{0,51,153},
    filecolor=[RGB]{0,51,153},
    allcolors=[RGB]{0,51,153}
}

%--------------------------------------------------
% SECTION NUMBERING
%--------------------------------------------------
\renewcommand{\thesection}{\arabic{section}.}
\renewcommand{\thesubsection}{\thesection\arabic{subsection}.}
\renewcommand{\thesubsubsection}{\thesubsection\arabic{subsubsection}.}

\renewcommand\thepart{}
\renewcommand\partname{}
\renewcommand\ptctitle{Table of Contents\\}

%--------------------------------------------------
% ICON SHORTCUTS
%--------------------------------------------------
\newcommand{\github}[1]{\href{#1}{\faGithubSquare}}
\newcommand{\youtube}[1]{\href{#1}{\faYoutubeSquare}}
\newcommand{\database}[1]{\href{#1}{\faDatabase}}
\newcommand{\gitlab}[1]{\href{#1}{\faGitlab}}
\newcommand{\globe}[1]{\href{#1}{\faGlobe}}


%--------------------------------------------------
% TEAM SUBMISSION ENVIRONMENT
%--------------------------------------------------
\newcounter{teamcount}

\newenvironment{teamsubmission}[2]{
    \clearpage
    \refstepcounter{section}
    \stepcounter{teamcount}
    \addcontentsline{toc}{section}{\theteamcount. #1}
    \begin{center}
        {\LARGE \bfseries #2 \par}
        \vspace{0.25em}
        {\large \textbf{Team:} #1 \par}
    \end{center}
}{
    \FloatBarrier
    \clearpage
}

%--------------------------------------------------
% AUTHORS & AFFILIATIONS
%--------------------------------------------------
% \newcommand{\authorsblock}[1]{%
%     \begin{center}
%         {\normalsize \textbf{#1}\par}
%     \end{center}
% }

% \newcommand{\affiliationsblock}[1]{%
%     \begin{center}
%         {\textit{#1}\par}
%     \end{center}
%     \vspace{0.5em}
% }
\newcommand{\authorsblock}[1]{}
\newcommand{\affiliationsblock}[1]{}

%--------------------------------------------------
% Roman numeral helper (I, II, III, ..., X, L, C, ...)
%--------------------------------------------------
\newcommand{\RNum}[1]{\uppercase\expandafter{\romannumeral #1\relax}}


\begin{document}

%--------------------------------------------------
% TITLE
%--------------------------------------------------
% \begin{center}
% \parbox{\textwidth}{%
% \centering
% \LARGE
% \setlength{\baselineskip}{1.1\baselineskip}
% \textbf{From Knowledge to Action: Outcomes of the 2025 Large Language Model (LLM) Hackathon for Applications in Materials Science and Chemistry}
% }
% \end{center}


% \vspace{0.25em}


% % ---------- Global Authors ----------
% \input{authors}

% \begin{center}
% \small
% Author names and ORCID identifiers are listed above; institutional affiliations appear in the respective team contributions. Institutional affiliations for authors not associated with a specific team are listed below:\\
% Andrew MacBride$^\dagger$ — University of California, Berkeley, California, US.\\
% Ben Blaiszik$^*$ — Argonne National Laboratory, Lemont, IL, USA.
% \end{center}
% \begin{center}
% \small
% $^*$Corresponding author: Ben Bliaszik (\href{blaiszik@uchicago.edu}{blaiszik@uchicago.edu}).\\$^\dagger$These authors also contributed substantially to compiling team results and other paper writing tasks.
% \end{center}

%--------------------------------------------------
% MAIN PAPER CONTENT
%--------------------------------------------------
\begin{tcolorbox}[
    enhanced,
    colback=mytheme,    % Background color
    colframe=mytheme,   % Match frame to background to hide border
    arc=3mm,            % Adjust this for more/less roundness
    boxrule=0pt,        % Explicitly remove border thickness
    left=15pt,          % Internal padding
    right=15pt,
    top=12pt,
    bottom=12pt
]
    \begin{minipage}{\linewidth}
    \centering
    {\par
    \fontsize{20pt}{22pt}\selectfont
    \fontseries{eb}\selectfont
    From Knowledge to Action: Outcomes of the 2025 Large Language Models (LLM) Hackathon for Applications in Materials Science and Chemistry
    \par}

    \vspace{15pt}

    {\large\textbf{
    2025 LLM Hackathon for Applications in Materials Science and Chemistry Participants$^1$}}\\
    {\normalsize $^1$A detailed contributor list can be found in the appendix of this paper.}
    
\end{minipage}

\vspace{10pt}
    In this article, we showcase 88 diverse applications of LLMs in materials science and chemistry as the outcomes from the third Large Language Model (LLM) Hackathon for Applications in Materials Science and Chemistry, engaging more than 350 participants across global hybrid locations. These diverse applications can be generally divided into two high-level categories: i) Knowledge Infrastructure (New ways to know things) and ii) Action Systems (New ways to do things). The appendix contains brief papers for each team submission, with accompanying links to code repositories, video demonstrations, and datasets where applicable. In addition to these team results, we provide a discussion of the hackathon event and its hybrid format, featuring 16 physical hubs distributed across four continents—North America, Europe, Asia, and Australia—complemented by a global online hub that enabled both local and virtual collaboration. Overall, the event revealed substantial advancements in LLM capabilities compared to last two years' hackathons, indicating ongoing expansion of LLM applications in materials science and chemistry research towards the realization of self-driving laboratories. The outcomes illustrate the dual utility of LLMs, serving both as multipurpose models for diverse machine learning tasks and as platforms enabling rapid prototyping of custom scientific applications, while also facilitating the deployment of fully automated workflows across computational and experimental environments. 
\end{tcolorbox}


\begin{bibunit}[ieeetr]
\section*{Introduction}
The integration of large language models (LLMs) into scientific workflows is rapidly reshaping how researchers perform knowledge discovery, automate routine tasks, and orchestrate end-to-end computational and experimental pipelines.
In chemistry and materials science, where progress is often constrained by heterogeneous data modalities, fragmented tooling, and rapidly evolving literature, LLMs have emerged as general-purpose interfaces that can bridge natural-language with domain knowledge, code, and laboratory procedures \cite{llm_science_1,llm_science_2,llm_matchem_1,llm_matchem_2}.
Early efforts have demonstrated promise across a wide spectrum of applications---from literature understanding and information extraction,
to code synthesis and simulation setup, to agentic systems that connect planning, execution, and reporting \cite{early_apps_1,early_apps_2}.

However, the pace of change in foundation models, tool ecosystems (\textit{e.g.}, retrieval-augmented generation and agentic frameworks), and deployment platforms make it difficult for individual research groups to continuously track capabilities, identify high-impact use cases, and validate what works in practice.
Community-driven, time-bounded innovation formats such as science hackathons offer a complementary mechanism: they have proven to be an effective tool for rapid prototyping in a friendly, competitive manner, while fostering interdisciplinary collaborations and innovations within the scientific community~\cite{pe2019understanding, nolte2020support, heller2023hack}.
Such hackathons can be organized cost-effectively while maximizing both their impact and reach through the use of social media, virtual platforms, and hybrid event structures.
Our previous hackathons on LLMs for Applications in Materials Science and Chemistry in 2023~\cite{jablonka202314} and 2024~\cite{llm2024} have already demonstrated the potential of LLM applications in diverse areas of materials science and chemistry, while also revealing persistent limitations and open challenges that motivate further work.

In this article, we discuss the results of the third LLM Hackathon for Applications in Materials Science and Chemistry, which engaged a large global community in a hybrid format.
The event produced a broad portfolio of submissions that collectively illustrate how LLMs are being used not only as multipurpose models for scientific machine-learning tasks, but also as enabling platforms for rapid prototyping of domain-specific applications.
We begin by detailing and analyzing the project submissions, organizing the results into two broad categories with in-depth sub-classifications, \emph{Knowledge Infrastructure} (new ways to know things) and \emph{Action Systems} (new ways to do things), reflecting the complementary roles of LLMs in (i) structuring, retrieving, and synthesizing scientific knowledge and (ii) executing or coordinating scientific work across computational and experimental environments.
We then briefly discuss the organisation of the 48-hour long event via one central virtual hub and 16 in-person sites, engaging more than 350 participants all over the world.
Finally, we conclude by suggesting future directions of LLM applications in materials science and chemistry, including increasingly autonomous and integrated research workflows aligned with the longer-term vision of self-driving laboratories \cite{self_driving_labs_1,self_driving_labs_2}.

\section*{Overview}
The third LLM Hackathon for Applications in Materials Science and Chemistry received 88 submissions from more than 350 participants across the world.
We broadly grouped submissions into Knowledge Infrastructure (\textit{New ways to know things}) and Action Systems (\textit{New ways to do things}) (Figure XX).
Knowledge Infrastructure describes submissions related to data curation and analysis, reproducible workflows, as well as educational and collaborative tools.
Action Systems are submissions describing new generative or inverse design models, and agentic systems that orchestrate computational chemistry tools.

\section*{Conclusion}
This is the main paper \textbf{Conclusion} content.

\section*{Contributors and Acknowledgements}
This paper is the result of contributions from all enthusiastic participants in the 2025 LLM Hackathon for Applications in Materials Science and Chemistry. Below, we list all contributors in alphabetical order by last name, apart from lead and corresponding authors. Only author names and ORCID identifiers are listed for all participants except the corresponding author. Individual authors can be contacted using their ORCID details or by contacting the corresponding author.
\vspace{5pt}

% ---------- Global Authors ----------
\input{authors}

\begin{center}
\small
$^1$ Data Science \& Learning Division, Argonne National Laboratory, Lemont, IL, USA.\\
$^2$ University of Chicago, Chicago, IL, USA\\
\vspace{5pt}
$^*$Corresponding author: Ben Bliaszik (\href{blaiszik@uchicago.edu}{blaiszik@uchicago.edu}).\\$^\dagger$These authors also contributed substantially to compiling team results and other paper writing tasks.
\end{center}


% \FloatBarrier
% \bibliographystyle{ieeetr}
% \bibliography{ref}
\putbib[ref]

\end{bibunit}

%--------------------------------------------------
% TABLE OF CONTENTS (AFTER MAIN PAPER)
%--------------------------------------------------
\clearpage
\tableofcontents
\newpage

%--------------------------------------------------
% TEAM SUBMISSIONS (ALL IN ONE FILE)
%--------------------------------------------------
\clearpage
\begin{bibunit}[ieeetr]

\input{teamsubmissions}

\putbib[submissionref]

\end{bibunit}

% \input{teamsubmissions}

% \FloatBarrier
% \bibliographystyle{ieeetr}
% \bibliography{submissionref}

\end{document}
